\section{What the midterm is}
\label{sec:midterm}

Thursday 2026-11-12, 18:10 to 19:40, at the Exam Centre, 255 McCaul Street. Ninety minutes, three questions, fifteen marks each. The instructions ask for three things in every answer: all the work shown, the assumptions stated, and units on every number. Constants you need are printed on the paper, and a crib sheet is handed out with it. Your own crib sheet and every other aid are not allowed.

\subsection*{What is examinable}

The course names sections of White, \emph{Fluid Mechanics}, 8th edition, as taught in lecture.

\begin{center}\small
\begin{tabular}{@{}l >{\raggedright\arraybackslash}p{6.2cm} l@{}}
\toprule
White & Content & This document \\
\midrule
3.1--3.2 & Basic laws, the Reynolds transport theorem & \secref{sec:rtt}{the Reynolds transport theorem} \\
3.3 & Conservation of mass & \secref{sec:mass}{conservation of mass} \\
3.4 & Linear momentum for a control volume & \secref{sec:momentum}{linear momentum} \\
3.5 & Frictionless flow, Bernoulli & \secref{sec:bernoulli}{Bernoulli} \\
3.6 & Angular momentum & \secref{sec:angular}{angular momentum} \\
3.7 & Energy & \secref{sec:energy}{the energy equation} \\
4.1 & The acceleration field & \secref{sec:acceleration}{the acceleration field} \\
4.2 & Differential continuity & \secref{sec:continuity}{differential continuity} \\
4.3 & Differential momentum, Navier--Stokes & \secref{sec:navierstokes}{Navier--Stokes} \\
4.6 & Boundary conditions & \secref{sec:bcs}{boundary conditions} \\
4.8 & Vorticity and irrotationality & \secref{sec:vorticity}{vorticity} \\
4.10 & Exact incompressible viscous flows & \secref{sec:exactflows}{exact viscous flows} \\
5.1--5.2 & Dimensional homogeneity & \secref{sec:homogeneity}{dimensional homogeneity} \\
5.3 & The Pi theorem & \secref{sec:pitheorem}{the Pi theorem} \\
5.4 & Nondimensionalizing the equations & \secref{sec:nondimensional}{nondimensional equations} \\
5.5 & Modelling and similarity & \secref{sec:similarity}{modelling and similarity} \\
\bottomrule
\end{tabular}
\end{center}

Not tested: 4.4 (differential angular momentum), 4.5 (differential energy), 4.7 (the stream function) and 4.9 (frictionless irrotational flow). They are in the course's Chapter 4 supplement and worth reading another time. The course says the best preparation is the lecture worked examples, the chapter summaries and the assignments. This document does not have those, so it sits beside them: it derives the theory, gives invented worked examples, and orders textbook problems into short ladders.

\subsection*{What the crib sheet gives, and what it does not}

The sheet handed out with the paper is two pages of formulas and a Moody chart, covering the whole course. For this midterm the useful parts are these.

\begin{itemize}[nosep,leftmargin=*]
\item Section 3 has the Reynolds transport theorem, relative velocity, mass, linear momentum, the force of a jet on a vane, angular momentum in two forms, and energy in head form with the kinetic energy factor $\alpha$.
\item Section 4 has the Pitot relations, which are Bernoulli in disguise.
\item Section 12 has the Pi theorem's count, the requirement that repeating variables be independent, and the standard dimensionless groups.
\item Section 13 has 2-D incompressible continuity, the vorticity and the element rotation, the strain rates, and the no-slip condition.
\item Section 14 has Euler's turbomachine head and torque, which is angular momentum for a pump or turbine.
\item Section 16 has the units and dimensions of the common quantities, and section 17 the properties of water by temperature.
\end{itemize}

Three gaps decide how to study.

\begin{enumerate}[nosep,leftmargin=*]
\item \textbf{No Navier--Stokes.} The sheet does not print the material derivative, continuity in three dimensions or in cylindrical coordinates, the Navier--Stokes equations in any form, or any exact viscous solution except the laminar pipe flow rate. The course announcement says the Navier--Stokes equations need not be memorized. Either the paper prints them, or the question gives what it needs. Plan for neither: each Chapter 4 section here shows how to write its equation from scratch in under two minutes, from Newton's second law and a box. That costs an hour of practice and removes the risk. Asking the course before Thursday which equations the paper prints is cheap.
\item \textbf{No assumptions.} Every equation on the sheet is printed without the conditions under which it holds. The marks for ``state your assumptions'' are exactly the part the sheet cannot give you. Each section here has a box, \emph{On the crib sheet}, that says what the sheet prints for that topic and what you must supply.
\item \textbf{No procedure.} The sheet has the Pi theorem's count, not the repeating-variable method, and the momentum equation, not the choice of control volume. Those two choices are the modelling calls of this midterm, and the next section is about them.
\end{enumerate}

\subsection*{How the marks are likely spread}

Three questions of fifteen marks in ninety minutes is thirty minutes a question, including the write-up. Questions of this length are usually multi-part, and the first parts carry the set-up marks: the control volume drawn, the assumptions listed, the right equation written down and reduced. A question whose final number is wrong but whose set-up is right usually keeps most of its marks. A question with the right number and no stated assumptions usually does not keep all of them. So the write-up order matters more here than in a problem set: sketch, assumptions, equation, reduction, numbers, units, check.

The three questions will probably follow the three chapters, one each: a control-volume question, a differential question, and a dimensional-analysis question. That is a guess from the announcement's chapter list, not something it says.
